\documentclass[11pt]{article}
\usepackage[T1]{fontenc}
\usepackage[letterpaper,margin=1.12in]{geometry}
\usepackage{amsmath,amssymb,amsthm}
\usepackage[hidelinks]{hyperref}
\newtheorem{theorem}{Theorem}[section]
\newtheorem{lemma}[theorem]{Lemma}
\newtheorem{proposition}[theorem]{Proposition}
\newtheorem{corollary}[theorem]{Corollary}
\theoremstyle{definition}
\newtheorem{definition}[theorem]{Definition}
\DeclareMathOperator{\Var}{Var}
\DeclareMathOperator{\sgn}{sgn}
\newcommand{\E}{\mathbb E}
\newcommand{\R}{\mathbb R}
\newcommand{\Vstar}{V^\star}
\newcommand{\cx}{\preceq_{\rm cx}}
\newcommand{\norm}[1]{\lVert #1\rVert_2}
\newcommand{\abs}[1]{\lvert #1\rvert}
\allowdisplaybreaks[1]
\clubpenalty=10000
\widowpenalty=10000
\setlength{\emergencystretch}{1em}
\hypersetup{
  pdftitle={Sharp Variance Comparisons on Laminar Information Intervals},
  pdfauthor={Jonas Weingart},
  pdfsubject={Technical manuscript, version 1.0 \textemdash{} 18 November 2026}
}
\title{Sharp Variance Comparisons for Conditionally Centered Scores on Laminar Information Intervals}
\author{Jonas Weingart}
\date{Technical manuscript, version 1.0 \textemdash{} 18 November 2026}

\begin{document}
\maketitle
\begin{abstract}
We maximize the variance of a deterministic weighted sum of bounded scores,
each revealed at a deterministic maturity and centered conditionally on its
issue-time information. A finite sign linear program retains the entire
matured-score history. For laminar schedules, we construct a single feasible
law that maximizes every convex payoff of the total. The construction
attains its bound with independently sampled child vectors even though
admissible later vectors may depend on earlier histories. Together with an
exact Gram recursion, it yields the sharp depth-two factor
$(1+1/\sqrt2)/2$ relative to a covariance semidefinite relaxation.
For one outer interval with at most two child stars of at most two active
leaves each, we prove analytically that the sharp factor is
$\gamma=3/(15-8\sqrt2)$. An exact computer-assisted proof extends this
factor to at most three such stars. Nonzero equality uses two active
stars and a unique weight pattern up to signs, relabeling and scale.
The unrestricted depth-three constant lies between $(1+4/\pi)/3$ and
$\gamma$ and remains undetermined.
\end{abstract}

\section{Introduction}

Overlap of issue-to-maturity intervals permits covariance, but issue-time
centering can constrain the joint law well beyond its second moments.
For a fixed schedule, we compare the largest attainable variance with
a covariance semidefinite program that fixes disjoint entries to zero.
The central question is when chronological conditioning allows the
worst law to be constructed one branch at a time.

Laminarity makes sibling intervals chronologically disjoint. Earlier
sibling scores therefore belong to every later sibling's incoming
information. The feasible kernels can depend on that information.
Our recursion works because the same symmetric child law dominates
every admissible conditional law in convex order. Maximization then
reduces to independent child vectors followed by a feasible root coupling.
The Gram problem has a parallel recursion based on orthogonal spans.
Combining these recursions exposes a trade-off between a Rademacher first
moment and a pairing-dependent Gram norm.

Convex comparison for bounded martingale differences and comonotone
maximization for specified marginals are established principles
\cite{bentkus,kaas,cheung,leskela}. The local fair-sign step below is
an instance of those principles. The argument that carries it through the
specified issue histories and attains it at all nodes is essential here.
Related lagged forecast inference uses conditional centering for different
objectives \cite{choe}. Graph Grothendieck problems compare deterministic
sign and unit-vector edge objectives \cite{laurent}. Our nonedge entries
are prescribed zeros and our sign variables obey full-history equations.
Sparse PSD decomposition addresses the matrix constraints
\cite{agler}. Gaussian rounding uses established PSD comparison techniques
\cite{nesterov,briet}. The sharp moment inequality and the sign chambers
also come from established work. The specific comparison proved here
optimizes the filtered-law numerator jointly with its Gram denominator.
The resulting recursion is distributional. Its statement carries no
polynomial-time bound for evaluating arbitrary weighted sums.

\section{Model and exact sign program}

\begin{definition}[Admissible scores]\label{def:model}
Fix a finite discrete clock and distinct deterministic supports
$I_i=(a_i,b_i]$ with $a_i<b_i$. On a filtered probability space,
a score $X_i$ satisfies
\[
 |X_i|\le 1,\qquad X_i\text{ is }\mathcal F_{b_i}\text{-measurable},
 \qquad \E[X_i\mid\mathcal F_{a_i}]=0.
\]
All scores maturing at a given time are available at issues at that time.
The weight vector $w$ is real and deterministic. A family is laminar
when any two supports are disjoint or one contains the other. Its depth
is the number of supports in a longest strict inclusion chain.
\end{definition}

The admissible class varies over filtered spaces and laws. All scores
have mean zero, so define
\begin{align}
 \Vstar &= \sup_{\rm admissible}\E\Big(\sum_i w_i X_i\Big)^2,\\
 U &= \max\{w^\top Cw : C=C^\top\succeq0,\ C_{ii}=1,\
           C_{ij}=0\text{ if } I_i\cap I_j=\varnothing\}.\label{eq:sdp}
\end{align}
Here $\Vstar$ is the worst variance for this law class. Replace $X_i$ by
$\operatorname{sgn}(w_i)X_i$ with either sign at a zero weight, and
conjugate $C$ by the same diagonal sign matrix. This preserves both
problems and permits nonnegative weights below. Other deterministic
score bounds can be absorbed into the weights. The identity is feasible
in~\eqref{eq:sdp}, so $U>0$ when $w\ne0$. All zero-weight instances
have both optima zero and require no ratio.

\begin{theorem}[Full matured-history sign LP]\label{thm:lp}
Let $M_i=\{j:b_j\le a_i\}$. The value $\Vstar$ equals the maximum of
$\sum_z p_z(w^\top z)^2$ over the following finite polytope:
\[
 p_z\ge0,\qquad \sum_{z\in\{-1,1\}^n}p_z=1,\qquad
 \sum_{z:z_{M_i}=h}p_z z_i=0
 \quad\bigl(i,\ h\in\{-1,1\}^{M_i}\bigr).
\]
For each convex function of the weighted total, maximizing its
expectation over this polytope gives the corresponding supremum.
\end{theorem}
\begin{proof}
On an independent product extension sample independent uniform
random variables $L_i$ on $[0,1]$. Release $L_i$ at $b_i$ and set
$Z_i=2\mathbf1_{\{L_i\le(1+X_i)/2\}}-1$.
Conditional on the original terminal field the signs are independent
with means $X_i$. Conditional Jensen therefore increases every convex
payoff, and the square increases in expectation by
$\sum_i w_i^2(1-\E X_i^2)$. At $a_i$ the variable $L_i$ is fresh.
Condition first on the original terminal field and the released uniforms
to integrate $L_i$, then project down to the original issue field
enlarged by the released uniforms. The conditional expectation of $X_i$
there equals $\E[X_i\mid\mathcal F_{a_i}]=0$ by independence. Projecting
down to the matured signs gives all the displayed equations.

Conversely, sample a sign vector with masses $p_z$ and at time $t$ take
$\sigma(Z_j:b_j\le t)$ as the filtration. Each positive-mass history
has conditional mean zero by its equation, while a zero-mass history
imposes no further condition. This realizes every feasible vector of
masses. Compactness of the nonempty polytope gives attainment.
\end{proof}

If $b_i\le a_j$, the tower property gives $\E X_iX_j=0$. Thus the sign
problem is bounded above by~\eqref{eq:sdp}. Pairwise zero covariances alone
do not enforce conditional centering. For three disjoint consecutive
intervals, let the signs be $R,S,RS$ for independent fair signs $R,S$.
The covariance matrix is the identity, but $RS$ is already determined
at its issue. In the opposite direction, admissible disjoint subtree
vectors may be dependent while each cross covariance vanishes.

\section{Laminar convex-order and Gram recursions}

For bounded random variables $Y,Y'$ write $Y\cx Y'$ when
$\E f(Y)\le \E f(Y')$ for every real convex $f$. Write $\sgn_0$ for
the sign with an independent fair tie coin at zero.

\begin{lemma}[Fair root extension]\label{lem:root}
If $S\cx S^\circ$, $S^\circ$ is symmetric, $Y$ is a fair sign with
arbitrary coupling to $S$, and $m\ge0$, then
\[
 S+mY\cx S^\circ+m\sgn_0(S^\circ).
\]
\end{lemma}
\begin{proof}
Fix convex $f$ and put $\lambda=(f(m)-f(-m))/2$ and
$g(s)=\max\{f(s+m)-\lambda,f(s-m)+\lambda\}$.
The convex function $g$ bounds $f(S+mY)-\lambda Y$ pointwise.
The expectations are bounded by $\E g(S^\circ)$ because $Y$ is fair
and $g$ is convex. The increment $f(s+m)-f(s-m)$ is nondecreasing
in $s$, and the two branches meet at zero. Hence $\sgn_0(S^\circ)$
selects a maximizing branch and is fair by symmetry. Substituting
this choice shows that $\E g(S^\circ)$ equals the stated upper payoff.
The argument also covers a zero weight and atoms at zero.
\end{proof}

The selected coupling is comonotone. More generally, with a fixed
bounded inner law $S$, the largest $\E[SY]$ over fair signs $Y$ is
$\min_t\E|S-t|$. For any $t$, fairness gives
$\E[SY]=\E[(S-t)Y]\le \E|S-t|$. At a median use the sign of $S-t$
off the median and choose the tie probabilities to make the sign fair.
If a root has a trivial incoming field and can observe the inner vector
at maturity, this coupling is admissible. The threshold zero in
Lemma~\ref{lem:root} is justified by symmetry of the reference law.

\begin{theorem}[Conditional laminar extremum]\label{thm:tree}
Order the children of every node $v$ of a finite laminar forest by time.
Independently sample their canonical vectors and a fresh fair tie coin
at $v$. Recursively define the canonical total
\[
 S_v=\sum_{c:\,c\text{ child of }v}T_c,\qquad
 T_v=S_v+w_v\sgn_0(S_v).
 \tag{*}\label{eq:tree}
\]
The empty sum is zero. For every admissible law and event
$E\in\mathcal F_{a_v}$ with $\Pr(E)>0$, the conditional distribution
of the weighted subtree total given $E$ is dominated by $T_v$ in
convex order. The canonical vector is admissible in its matured-sign
filtration. Independent maximal components give a convex-order
maximum for the forest total and
\[
 \Vstar=\sum_{v\text{ maximal}}\E T_v^2.
\]
\end{theorem}
\begin{proof}
Prove the eventwise statement by induction on the subtree. For bounded
scores, endpoint rounding as in Theorem~\ref{thm:lp} also applies
under the law conditioned on $E$. In particular, $E$ belongs to the
original terminal field, so conditional Jensen given that field and
averaging on $E$ yield the convex comparison with the rounded signs.
It suffices to prove the induction for signs. Conditioning on $E$
preserves every descendant's issue constraint because $E$ belongs to
each descendant's issue field. For any event $F$ in that issue field,
the original centering equation integrated on $E\cap F$ gives the
conditional-law equation on $F$.
A leaf is fair under this conditioning and has the canonical law.

Under the law conditioned on $E$, fix the complete sign vectors of
the children preceding the last child. Their maturities are no later
than the last child's issue. Each positive-mass joint history, together
with $E$, is an incoming event for that subtree. On this history the
earlier total is a fixed number. The induction hypothesis applied to
the convex payoff translated by that total bounds the expected payoff
by averaging against the same canonical last-child law for every
history. Such averaging leaves a convex function of the earlier sum.
Iterating backwards over the children replaces the adaptive joint sum
by the sum of independent canonical child totals in convex order.
This reference sum is symmetric. The actual root sign is fair on $E$.
Lemma~\ref{lem:root} proves the eventwise comparison with~\eqref{eq:tree}.

For attainment generate the canonical signs from independent primitive
tie coins at the nodes. A node's sign is a function of its own subtree's
coins and is fair by symmetry, including at zero. At its issue every
descendant is still immature because each has positive length and starts
at or after that issue. Every ancestor is also immature. The matured
scores are in separated branches and are functions of coins
outside the node's subtree. Thus the sign is conditionally fair at issue. The
matured-sign filtration reveals all coincident maturities together and
makes the assigned signs measurable at maturity. This establishes
simultaneous feasibility. Apply the same backwards argument to the
ordered maximal components, without an additional root, to obtain the
forest convex-order maximum. Independent component totals are centered,
so the square yields the formula for $\Vstar$.
\end{proof}

Independent child vectors are one attaining choice. For example, with
two disjoint two-leaf stars of unit weights inside one root, the later
root may copy its first leaf if the earlier first leaf is positive and
its second leaf otherwise. Its conditional total has the same law on
both histories, while its joint vector depends on the first star.
The outer sign can still attain the optimum.

\begin{proposition}[Gram recursion]\label{prop:gram}
Define $u_v=w_v+\sqrt{\sum_{c:\,c\text{ child of }v}u_c^2}$. Then
$U=\sum_{v\text{ maximal}}u_v^2$.
\end{proposition}
\begin{proof}
Represent $C$ by unit Gram vectors. The spans of distinct child
subtrees are orthogonal because their supports are disjoint. The
largest possible child-sum norm is bounded by the square root of the
sum of their squared optima. Adding the parent increases it by at most
$w_v$. To attain both bounds, embed optimal child systems in orthogonal
spaces and align the parent with their weighted sum. If it vanishes,
use a fresh orthogonal unit vector. No disjointness zero inside the
subtree involves its parent. At a maximal component the same construction
leaves its span orthogonal to every other maximal component.
\end{proof}

\section{Sharp comparison at depth two}

A star consists of a root with its pairwise disjoint leaves. For one
such component write $r$ for the root weight, $b=\norm{q}$ for the leaf
norm and $A=\E|\sum_i q_i\epsilon_i|$ for the first absolute moment of
the independent fair leaf signs. The recursions give
\[
 \Vstar=r^2+b^2+2rA,\qquad U=(r+b)^2.
\]

\begin{theorem}[Depth two]\label{thm:depth2}
For every finite laminar forest of depth at most two,
$\Vstar\ge c_2U$ with sharp constant $c_2=(1+1/\sqrt2)/2$.
If $U>0$, equality holds precisely when each positive-weight
component has exactly two equal nonzero leaf magnitudes and root
magnitude equal to its leaf norm.
\end{theorem}
\begin{proof}
Subtract $c_2U$ on each component:
\[
 \Vstar-c_2U=(1-c_2)(r-b)^2+2r(A-b/\sqrt2).
\]
The real sharp $L^1$ Khintchine inequality gives $A\ge b/\sqrt2$
\cite{szarek,haagerup}. For $b>0$, equality in that inequality requires
exactly two equal nonzero magnitudes \cite{de}. A positive-weight
component with zero deficit therefore has $r=b>0$ and the described
leaves. The two forest objectives add across maximal components.
For sharpness take the root $(0,2]$, leaves $(0,1],(1,2]$, and
weights $(\sqrt2,1,1)$. Then $A=1$, $\Vstar=4+2\sqrt2$, and $U=8$.
At $r=b>0$ the comparison reduces to the sharp Khintchine bound.
\end{proof}

\section{An outer interval with two or three small stars}

Consider one outer root of magnitude $r$ whose children are stars
with at most two nonzero leaf weights each. For a child with root
magnitude $m$ and leaf magnitudes $p\ge q\ge0$, the canonical total in
Theorem~\ref{thm:tree} has law $(m+p)\epsilon+q\eta$ for independent fair
signs. For $p>q$ the canonical root copies the heavier sign. For $p=q$
copying that sign gives the same distribution as the fair tie rule.
When $p=q>0$, the two cancelling leaf configurations are equiprobable,
and the copied sign takes opposite values on them. If $p=q=0$,
a fresh fair sign gives total $m\epsilon$. Zero or absent leaves may
be padded algebraically. The corresponding child Gram norm is
$m+\sqrt{p^2+q^2}$.

At fixed effective pair $(a,b)=(m+p,q)$ we have $a\ge p\ge b\ge0$.
The Gram norm $a-p+\sqrt{p^2+b^2}$ is largest at $p=b$, strictly so
for $b>0$ and $p>b$. Thus the replacement $(m,p,q)=(a-b,b,b)$
preserves the canonical inner law and weakly increases $U$. Set
$k=\sqrt2-1$. In the reduced family, independently over the effective
pairs $(a_j,b_j)$, put
\begin{align*}
 S&=\sum_j(a_j\epsilon_j+b_j\eta_j), & A&=\E|S|,\\
 Q&=\sum_j(a_j^2+b_j^2), & B&=\sqrt{\sum_j(a_j+kb_j)^2}.
\end{align*}
The quotient is
\begin{equation}
 \frac{\Vstar}{U}=\frac{r^2+2rA+Q}{(r+B)^2}.
 \label{eq:ratio}
\end{equation}
A reduced bound applies to the original problem because its numerator is
unchanged and its denominator can only increase. For six unit inner
weights and $r=1$, the exact values are $A=5/2$, $Q=10$ and
$B=2+\sqrt2$. Thus the quotient is $16/(11+6\sqrt2)<c_2$.

\begin{lemma}[Dominant effective coefficient]\label{lem:dominant}
If the largest effective coefficient is at least the sum of all others
and $U>0$, the quotient in~\eqref{eq:ratio} is strictly above
$\gamma=3/(15-8\sqrt2)$.
\end{lemma}
\begin{proof}
For a nonzero effective vector let $a,b$ be its two largest coefficients,
padding zeros when needed. Conditioning on the other signs gives $A=a$.
To upper-bound $B$ we may pair $a$ with $b$. Indeed, if their partners
were $c,d$, replacing $(a,c),(b,d)$ by $(a,b),(c,d)$ increases the
squared norm by $2k(b-c)(a-b-c-d)\ge0$ before sorting the last pair.
Sorting that pair can only increase it. Let $L$ sum the coefficients
other than $a,b$, $e=a+kb$, and $J$ be the remaining pair norm. Since
$(x+ky)^2\le x(x+y)$ for $x\ge y\ge0$, $J^2\le bL$ and
$e\ge(1+k)b+L$. If $J>0$,
\[
 \frac{J}{\sqrt{e^2+J^2}+e}
 \le \frac{\sqrt{bL}}{2((1+k)b+L)}
 \le \theta:=\frac1{4\sqrt{1+k}} .
\]
Thus $B\le a+kb+\theta J$, also for $J=0$.
Cauchy--Schwarz gives $J^2\le(1+k^2)Q_{\rm rest}$ for the remaining
squared coefficients $Q_{\rm rest}$. Apply it again to
$(r+a,b,\sqrt{Q_{\rm rest}})$ to bound~\eqref{eq:ratio} below by
\[
 \frac1{(1+k^2)(1+\theta^2)}
 =\frac{8\sqrt2+8}{16\sqrt2+1}>\gamma .
\]
The strict comparison follows from $8\sqrt2>11$. If the effective vector
is zero, the nonzero root alone gives quotient one.
\end{proof}

\begin{theorem}[Seven weights]\label{thm:seven}
For one outer root and at most two disjoint child stars with at most
two active leaves each, $\Vstar/U\ge\gamma$ whenever $U>0$, where
\[
 \gamma=\frac3{15-8\sqrt2}.
\]
This is the minimum over the class of such trees and weights. Equality
requires two child stars with two leaves available in each. Up to
common positive scaling and relabeling, the equality magnitudes are
outer weight $2$, both child root weights $(3\sqrt2-1)/2$, and four
unit leaves. Every other leaf has weight zero. Signs are arbitrary.
\end{theorem}
\begin{proof}
Sort the four effective coefficients $a\ge b\ge c\ge d\ge0$.
The numerator is invariant under pairings. For two pairs denote the
squared norm by $B_{ab,cd}^2=(a+kb)^2+(c+kd)^2$ and use analogous
notation for other pairings. Then
\begin{align*}
 B_{ab,cd}^2-B_{ac,bd}^2&=2k(b-c)(a-b-c-d),\\
 B_{ac,bd}^2-B_{ad,bc}^2&=2k(a-b)(c-d).
\end{align*}
The dominant case, including its boundary, is strict by
Lemma~\ref{lem:dominant}. In the other case a largest norm uses the
cross pairing $(a,c),(b,d)$. The moment formula is
\[
 A=\max\{a,(a+b+c)/2,(3a+b+c+d)/4\}.
\]
Indeed, averaging the sign of $a$ yields the maximum of $a$ and the
absolute remaining sum. Only $b+c+d$ and $b+c-d$ among the four
possible remaining magnitudes can exceed $a$.

Write $(a,b,c,d)=(x+u,x-u,y+v,y-v)$ in the nondominant case.
Then $0\le u,v\le y$ and $x\ge y+u+v$. For the reference
$(x,x,y,y)$ paired as $(x,y),(x,y)$, write
$A_0=x+y/2$, $B_0=\sqrt2(x+ky)$, $Q_0=2(x^2+y^2)$.
Put $d_0=\max(u,v)$. Evaluating the second and third branches of
the moment formula gives $A-A_0\ge\max(v/2,u/2)=d_0/2$.
The squared differences are $B^2-B_0^2=2(u+kv)^2$ and
$Q-Q_0=2(u^2+v^2)\ge c_2(B^2-B_0^2)$.
For $d_0>0$ rationalization gives
\[
 B-B_0\le \frac{(u+kv)^2}{\sqrt2(x+ky)}
 \le d_0\frac{u+kv}{\sqrt2d_0+u+v}\le k d_0.
\]
The second step uses $u+kv\le\sqrt2d_0$ and
$x+ky\ge\sqrt2d_0+u+v$. The last step uses $(1-k)u\le k\sqrt2d_0$. In the deficit
$r^2+2rA+Q-\gamma(r+B)^2$, the linear term's increase is
nonnegative because $\gamma k<1/2$, and the constant term strictly
increases when $d_0>0$ because $c_2>\gamma$. When $d_0=0$ the
reference is the original sorted vector. It suffices to optimize
the reference and any equality must be symmetric.

Put $s=\sqrt2$, $z=(r,x,y)^\top$ and $\ell=(1,s,2-s)^\top$.
The reference quotient is $z^\top M z/(\ell^\top z)^2$, where
\[
 M=\begin{pmatrix}1&1&1/2\\1&2&0\\1/2&0&2\end{pmatrix}.
\]
The Schur complement of the lower block is $3/8>0$.
Cauchy--Schwarz in the $M$ inner product yields the reciprocal of
$\ell^\top M^{-1}\ell=(15-8s)/3$ as a lower bound. The vector
\[
 M^{-1}\ell=\bigl((4-2s)/3,\ (5s-4)/6,\ (2-s)/3\bigr)^\top
\]
has $r>0$ and $x>y>0$, so its positive ray attains this bound. At
$y=1$ it gives $r=2$ and $x=x_0:=(1+3s)/2$. Both maximizing
pairs must be $(x_0,1)$ because pairing the like entries strictly
reduces $B$. Positive lighter coefficients force equal leaves
on reversing the within-star replacement. The formula includes
zero padding, which cannot give equality on a missing required
leaf or star. At the witness
$\Vstar=(39+18\sqrt2)/2$ and $U=(99-14\sqrt2)/2$.
\end{proof}

Fang and Wang \cite[Theorem 1]{fangwang} prove that, for sorted
nonnegative unit vectors of dimension at least four at distance
$0\le\tau\le\sqrt{2-\sqrt2}$ from $(1,1,0,\ldots)/\sqrt2$,
a profile $(x,x,y,y,0,\ldots)$ uniquely minimizes the first absolute
Rademacher moment. The normalized equality vector in
Theorem~\ref{thm:seven} has distance about $0.36243$, within that
range. The quotient optimizes over the distance, outer weight and
matching as well. For the same coefficient profile, crossing
heavy with light gives $B^2=2(x_0+k)^2$, while pairing like entries
gives $B^2=Q$. These pairings share the first moment and distance.

\begin{theorem}[Three small child stars]\label{thm:three}
With at most three disjoint child stars and at most two active leaves
per star, the sharp factor in $\Vstar\ge\gamma U$ is the same
$\gamma$. At $U>0$, equality uses the weights of
Theorem~\ref{thm:seven} with any remaining star of weight zero at every
node. If all three stars have positive total weight, the bound is
strict. The infimum over that subclass is $\gamma$, approached on
trees with at least two child stars that each have two leaves
available.
\end{theorem}
\begin{proof}
Minimize the quadratic deficit at fixed coefficients over $r\ge0$.
Its minimizer is $r=(\gamma B-A)_+/(1-\gamma)$, so the target is
equivalent to
\begin{equation}
 (\gamma B-A)_+^2\le(1-\gamma)(Q-\gamma B^2).
 \label{eq:deficit}
\end{equation}
Here $Q\ge c_2B^2$ and $A\le\sqrt Q\le B$. The last comparison
follows pairwise from $(a+kb)^2-a^2-b^2=2kb(a-b)$. For
$H=(1-\gamma)/\gamma$ we have
\[
 (\gamma B-A)^2-(1-\gamma)(Q-\gamma B^2)
 =\gamma\bigl((B-A)^2-H(Q-A^2)\bigr).
\]
Consequently the norm inequality
\begin{equation}
 B-A\le\sqrt{H(Q-A^2)}
 \label{eq:norm-test}
\end{equation}
suffices. A strict norm inequality excludes equality.
If $\gamma B\le A$ and $B>0$, strictness follows already from
$Q>\gamma B^2$. A nonzero outer weight with $B=0$ gives ratio one.
Pad to three pairs with zero pairs where needed. Lemma~\ref{lem:dominant}
handles the dominant region. For all other coefficients,
Appendix~\ref{app:tape} proves~\eqref{eq:norm-test} using a finite
exact certificate. Its cone inequalities hold on
all nonnegative generator combinations.

The appendix also shows that equality in the nonzero case forces a
positive multiple of $(x_0,x_0,1,1,0,0)$ in sorted coefficients.
Only two heavy--unit pairs and a zero pair maximize $B$ on this
ray. The deficit strictly decreases with $B>0$ at fixed $r,A,Q$,
so all other matchings give strictness. The unique minimizing
outer weight is $2$ in units of the lighter active coefficient,
and the equality leaves follow from the within-star replacement.
The seven-weight witness attains equality. Adding positive weights
to the third star and taking them to zero gives the same infimum
by continuity of the finite formulas.
\end{proof}

\section{Further bounds and limitations}

\begin{proposition}[Unrestricted depth three]\label{prop:c3}
For the infimum $c_3$ of $\Vstar/U$ over laminar forests of depth at
most three with $U>0$,
\[
 \frac{1+4/\pi}{3}\le c_3\le\gamma .
\]
\end{proposition}
\begin{proof}
After taking weight magnitudes the Gram recursion admits a maximizing
system with nonnegative coordinates: place child systems on disjoint
coordinate sets and align each root with the nonnegative sum, using
a fresh coordinate at zero. Let $C$ be this covariance matrix and let
$D=\sum_iw_i^2$. For a Gaussian vector with covariance $C$, each
coordinate is independent of the joint vector of coordinates matured
at its issue because its cross-covariances with this vector vanish.
Gaussian signs therefore satisfy the history LP. The Gaussian sign
correlation formula gives $(2/\pi)\arcsin C_{ij}$. Since $C_{ij}\ge0$,
the resulting objective is at least $D+(2/\pi)(U-D)$.

For a subtree of depth $d$ let $D_v$ be its sum of squared weights.
Inductively $u_v^2\le dD_v$. At a nonleaf use
$\sum_c u_c^2\le(d-1)(D_v-w_v^2)$ and apply
Cauchy--Schwarz to the two terms of the Gram recursion. This gives
$U\le3D$ in depth three, and the lower bound follows. The upper bound
uses the attaining weights of Theorem~\ref{thm:seven}.
\end{proof}

The equality $c_3=\gamma$ is open. The three-star theorem restricts
both the number of child stars and the number of active leaves.
The depth-three formula still follows from the general recursions
with either restriction removed, but its first moment and Gram norm
then need a new joint comparison.

\subsection{Four-star partial results}

The established four-star analysis covers a relative neighbourhood of
the three-star boundary, four identical effective pairs, and the slice
with two zero lighter coefficients. Its explicit local estimate is useful
for identifying an obstruction to a naive deletion argument. Choose two
pairs $(x+u,y+v),(x-u,y-v)$ and let $Z$ be the sum of the remaining
pairs, with Gram norm $B_Z$. Set
\[
 B_0=\sqrt2(x+ky),\quad D_0=2(u+kv)^2+B_Z^2,\quad
 \sigma=\tfrac14\sqrt{c_2/2},\quad \lambda=\sigma/\gamma .
\]
Here $B^2=B_0^2+D_0$. Relative to the two identical pairs $(x,y)$,
the increase in the first moment is at least $\sigma\sqrt{D_0}$.
To see this, condition on the cancellation of the two large signs and
the two small signs. On the joint cancellation event of probability
$1/4$, use sharp Khintchine on $2u\epsilon+2v\eta+Z$.
On the other events Jensen bounds the contribution below by the
corresponding reference contribution. The second moment increases by
at least $c_2D_0$ by Cauchy--Schwarz. Consequently
$\sqrt{D_0}/B_0\le 2\lambda/(1-\lambda^2)$ for $B_0>0$ gives
$A-A_0\ge\gamma(B-B_0)$, and the deficit at $\gamma$ cannot decrease.
On the normalized four-pair domain $B=1$, the minimized deficit is
continuous. On the boundary with a zero full pair it is positive
except at the equality rays of Theorem~\ref{thm:three}, where the
explicit estimate applies. Compactness yields a uniform neighbourhood
of that boundary with an existential radius.

The identical-pair slice has exact minimum
$367/(1919-1040\sqrt2)>\gamma$ in the frozen four-star analysis.
For four copies of $(a,b)$, the magnitude of the sum of four
independent fair signs has masses $3/8,1/2,1/8$ at $0,2,4$.
Using two independent copies gives
$A=3a/2+9b/16+(2b-a)_+/8$, $Q=4(a^2+b^2)$ and $B=2(a+kb)$.
Discarding the positive-part term gives a quadratic quotient
with numerator matrix and denominator vector
\[
 M_4=\begin{pmatrix}1&3/2&9/16\\3/2&4&0\\9/16&0&4\end{pmatrix},
 \qquad \ell_4=(1,2,2k)^\top .
\]
Its Schur complement is $367/1024>0$. The lower bound is
$(\ell_4^\top M_4^{-1}\ell_4)^{-1}=367/(1919-1040\sqrt2)$.
Direct substitution gives positive coordinates for $M_4^{-1}\ell_4$
with $a>2b$, so the discarded term vanishes at the minimizing ray.
Two singleton pairs $(a,0),(b,0)$ may be merged into one sorted pair
without changing the Rademacher law or decreasing the Gram norm.
Theorem~\ref{thm:three} then applies. General deletion is obstructed.
For instance, appending $(1/2,1/4)$ to $(5,2),(4,2),(2,0)$ leaves
$A$ unchanged because the old sum is odd and the added magnitude
is at most $3/4$. At the old minimizing outer weight the increase
in the Gram denominator is sufficient to decrease the quotient.
The companion source map identifies the exact comparisons in S08
and the reproduced fixture. The resulting quotient remains above
$\gamma$.

Selected closed four-star chambers have separate exact certificates.
The sharp child of chamber C illustrates a distinction between the
variance inequality and its certificate. In the fixed-projection
quadratic comparison a zero on a two-ray support imposes a kernel
condition on the PSD part of the aggregate residual. The certified
residual at this split and parameter $\delta=56/25$ requires rank two
in a PSD-plus-entrywise-nonnegative decomposition. The other child has
a certificate of rank one after the projection is changed. Their union
covers C. This is a local result for specified geometric data.

Copositive zero supports, SPN decompositions and simplicial verification
are established machinery \cite{hildebrand,shaked,bundfuss,bundfussadaptive}. A negative
value of one aggregate on a nonnegative vector excludes an SPN
certificate for that aggregate at any rank. The majorization can be
changed while retaining the original variance problem. For the
polynomial comparisons used with shifted projections, the Bernstein
basis is nonnegative and sums to one on the closed interval
\cite{funfzig}. Nonnegative coefficients give a nonnegative polynomial.
Strict positivity of an anchor requires a strictly positive lower
bound on the whole interval. All strictly positive coefficients give
such a bound, as in the two frozen positive examples. For instance,
the polynomial $t$ has nonnegative coefficients on $[0,1]$ and
vanishes at the endpoint. The full four-star inequality remains open.

\subsection{Interpretation and limitations}

The sharp comparisons concern two optima. Every admissible process
satisfies $\Var(\sum_i w_iX_i)\le\Vstar\le U$, and the identically
zero scores show that a positive multiple of $U$ gives no variance
floor for an individual process. A prescribed outcome mechanism or
richer fixed filtration can shrink the attainable class. Adaptive
weights, random schedules, and intervals that cross without nesting
fall outside the stated sharp laminar theorems. Marginal centering
alone cannot replace issue-time centering. These results supply no
forecast accuracy, sequential-testing or profitability conclusion.

\appendix
\section{Exact certificate for the three-star bound}\label{app:tape}

The certificate establishes~\eqref{eq:norm-test} on the
nondominant part of the sorted nonnegative cone in six coordinates.
It covers each of the fifteen perfect matchings. Homogeneity permits a
unit-sum slice for ray computations. The numerical search that proposed
the subdivision and active sets supplies no inequality or coverage step.
The input to exact verification is the saved finite plan and the explicit
geometry below.

\subsection{Polyhedral partition}

Sort $x_1\ge\cdots\ge x_6\ge0$ and retain all index matchings,
including when some values tie or vanish. For a matching
$\mathcal M$ let $P_\mathcal M x$ consist of $x_i+kx_j$ for each
pair $i<j$. Then $B=\norm{P_\mathcal M x}$ and
\[
 A(x)=\frac1{32}\sum_{s\in\{1\}\times\{-1,1\}^5}\abs{s^\top x}.
\]
Hausmann and Rodriguez \cite[Table (1--1), Sections 4--5]{hausmann}
classify the same length arrangement. Each wall equates the lengths
in a subset with those in its complement. Writing their difference
as a sign form, reversing the sorting convention and taking closures
identifies the walls here. Their chamber counts in six and eight
coordinates are 21 and 2,470, respectively. The exact verification
below reconstructs the six-coordinate geometry and uses it for a
matching-dependent inequality.

In gap variables $d_i=x_i-x_{i+1}$ with $x_7=0$, the coefficients of
$s^\top x$ are the successive prefix sums of $s$. Exactly twelve
forms have both positive and negative gap coefficients and need
both orientations. Every other form has a constant weak sign.
For each of the $2^{12}$ weak orientations, the intersection with
the nonnegative gap cone is pointed. Every extreme ray has five
linearly independent active homogeneous normals, with defining
equalities also available on lower-dimensional chambers.
Enumerate all five-normal choices from the six gap facets and
twelve walls and take their signed cofactors. Retain nonzero
nonnegative null vectors, convert gaps to coefficients, and normalize
each to primitive integer coordinates. For each orientation retain
all rays satisfying its weak inequalities. This procedure lists
every extreme ray and hence generates every closed cone.
Dropping duplicate ray sets and proper subsets preserves the union.
The exact result is eighteen rays and twenty-one maximal chambers.
Each maximal cone has six independent generators, checked by its
determinant. The one cone with strict dominance in its interior is
handled by Lemma~\ref{lem:dominant}. Its shared faces remain covered.
Zero weights and ties belong to closed chambers, and the origin
gives zero inner moments.

\subsection{A certificate on one closed cone}

Throughout one chamber $A=\ell^\top x$ with $\ell$ a signed
average of sign vectors. Then $|\ell^\top z|\le A(z)\le\norm z$
in the entire ambient space. Put $N=I-\ell\ell^\top$. It is positive
definite on each nondominant maximal chamber. Indeed,
$\norm\ell=1$ would give $A(\ell)=1$, so the absolute Rademacher
sum at $\ell$ would be constant by Cauchy--Schwarz. The degree-two
sign coefficients of its square imply $\ell_i\ell_j=0$ for $i\ne j$.
Positivity and $A(x)\ge x_1$ on a full-dimensional sorted chamber
then force $\ell=e_1$. Averaging first over the largest sign gives
$A(x)=x_1$ exactly in the dominant region.

For generators $v_i$ normalized to unit coordinate sum put
$G_{ij}=v_i^\top N v_j$, and choose
$\beta_i\ge \norm{P_\mathcal M v_i}-A(v_i)$. For a certificate
vector $h$ verify
\begin{equation}
 Gh\ge\beta,\qquad h^\top Gh\le H,\quad H=4-8\sqrt2/3 .
 \label{eq:check}
\end{equation}
With $x=\sum_i t_i v_i$ for $t_i\ge0$ and $y=\sum_i h_i v_i$,
the triangle inequality,~\eqref{eq:check} and Cauchy--Schwarz
in $N$ give
\[
 B(x)-A(x)\le\sum_i t_i\beta_i
 \le x^\top Ny\le \sqrt{x^\top Nx}\sqrt{h^\top Gh}.
\]
This is the required continuous inequality. The entries of $h$
may have either sign. At a stored active set $J$ the checker sets
$h_{J^c}=0$, solves $G_{JJ}h_J=\beta_J$ and checks every row of
$Gh\ge\beta$. The norm equals $h_J^\top\beta_J$.

A ray norm is exact for zero, one nonzero entry or two equal nonzero
entries of $P_\mathcal M v_i$. Otherwise the least upper bound on
the grid of denominator $10^{12}$ is used, obtained by comparing
squares. The ray norm is at most one by unit-sum normalization.
All inequalities use $\mathbb Q(\sqrt2)$. For opposite signs of
the two rational parts of $a+b\sqrt2$, testing its sign reduces
to comparing $a^2$ and $2b^2$ and selecting the dominant sign.
The linear solves have rational Gram matrices and use positive
pivots. All field comparisons are exact.

\subsection{Subdivision plan and verification interface}

Sort the primitive integer rays lexicographically. A chamber is
identified by its increasing tuple of ray indices in that list.
Visit chambers lexicographically and omit the dominant chamber.
Within each chamber visit matchings in the recursive order that
pairs the least unused index with each remaining index in increasing
order and recurses on the others. This yields $20\cdot15=300$
root cases. A base64 block in the preserved S07 document decodes
to the preorder traversal of those binary trees.
Bytes $1$ through $63$ are nonempty active-set masks on the six
current generators. Bytes $64$ through $78$ select one of the
fifteen pairs of generator slots in lexicographic order.

For a split on slots $i<j$, replace $v_j$ by $(v_i+v_j)/2$ in
the first child and $v_i$ by the same midpoint in the second.
Both children remain simplicial. If $t_i\ge t_j$, use $2t_j$
on the midpoint and $t_i-t_j$ on $v_i$ in the first child.
If $t_i\le t_j$ use the reversed representation in the second.
They cover the entire parent and share the equality face. The
verifier visits both branches and consumes the plan exactly.
It checks 1,854 nodes, including 777 splits and 1,077 terminal
certificates, at a maximum depth of six. The node count also satisfies $1854=2(1077)-300$.

On a strict norm leaf,~\eqref{eq:check} gives a strict bound for
all nonzero $x$ by positive definiteness of $N$. Exactly 34 leaves
saturate the norm condition. For every such leaf the verifier
checks that $y$ is a positive multiple of $(x_0,x_0,1,1,0,0)$.
Equality in Cauchy--Schwarz forces $x$ onto this ray. This supplies
a necessary equality direction on every terminal leaf. Sharpness
itself follows from the explicit attaining weights in the main
proof.

The exact plan, embedded in the S07 copy, and the two Python
verifiers accompany the manuscript under \path{certificate/}.
See \path{README.md} for the two replay commands. S07 enumerates
rays by cofactors and bounds the norm by field comparisons. S08
independently reconstructs the partition by clipping the sorted
simplex against the walls, including all opposite-side intersections
and filtering by active-normal rank. It uses full sign sums,
an alternative field representation, integer square-root brackets
and rational $LDL^\top$ solves. The two implementations verify the
same proposed subdivision and active-set data. S08's optional
eight-coordinate geometry run counts 2,470 maximal cells, 206
rays and a maximum of 37 rays in a cell. Its checks concern these
partition statistics, with shared faces included by the geometry.
A certificate of all four-star inequalities would require separate
checks on the remaining domain.


\phantomsection
\pdfbookmark[1]{Declaration of AI-Assisted Research}{ai-declaration}
\section*{Declaration of AI-Assisted Research}

This research was conducted under human direction with substantial assistance from generative artificial intelligence systems.

AI systems contributed to mathematical exploration, conjecture formulation, the development and refinement of proofs, computational experimentation, the construction and verification of exact mathematical certificates, literature investigation, and manuscript preparation.

The research underwent iterative mathematical and computational scrutiny, including exact-arithmetic verification and critical assessments by separate AI models. These assessments do not constitute independent human peer review or formal proof-assistant verification.

The human author directed the research programme, evaluated and selected the results presented, and assumes responsibility for the manuscript's mathematical claims, source attribution, methodological transparency, and limitations.

The mathematical arguments and computational certificates are provided to enable independent examination and reproduction.

\bibliographystyle{plain}
\begingroup
\raggedright
\bibliography{references}
\endgroup
\end{document}
